\documentclass[10pt,a4paper]{amsart}
\usepackage[margin=19mm]{geometry}
\usepackage{amsmath,amssymb,mathtools}
\usepackage[scr=rsfs,cal=euler]{mathalfa}
\usepackage{xfrac}
\usepackage[unicode,hidelinks,bookmarks=false]{hyperref}
\newcommand{\ie}{i.e.\ }
\newcommand{\cf}{cf.\ }
\newcommand{\resp}{respectively\ }
\newcommand{\Subsection}{\subsection}
\providecommand{\nobreakdash}{\nobreak\hbox{-}}
\setlength{\emergencystretch}{2em}
\multlinegap0pt
\usepackage{mathtools}
\usepackage{amssymb}
\newtheorem{lemma}{Lemma}
\newtheorem{proposition}{Proposition}
\DeclareMathOperator{\Hom}{Hom}
\DeclareMathOperator{\Mat}{Mat}
\DeclareMathOperator{\rk}{rk}
\newcommand{\vF}{\mathbb F}
\title{A General Construction of Codes from Drinfeld Modules}
\author{Alessandro Giannoni, Giacomo Micheli, Mihran Papikian}
\date{Source: arXiv:2609.01484v1 (CC BY 4.0)}
\begin{document}
\maketitle

\noindent\emph{Source and licence.} A General Construction of Codes from Drinfeld Modules. Version arXiv:2609.01484v1 (CC BY 4.0), distributed by arXiv under the Creative Commons Attribution 4.0 International licence, http://creativecommons.org/licenses/by/4.0/, as recorded in the arXiv licence metadata for this exact version and shown on the abstract page https://arxiv.org/abs/2609.01484v1. Full licence notice: \texttt{licenses/arxiv-2609.01484-CC-BY-4.0.txt}.\par\smallskip

\noindent\textit{Editorial scope.} Counted unit: the article's proposition proposition (source0.tex lines 937--947). The article's complete original proof of it is delivered with it (source0.tex lines 949--1033, one \texttt{proof} environment of 85 lines). No mathematics is written, repaired or re-derived: the statement and the proof are the article's own lines, in the article's own notation, and no retained line is substituted or re-flowed.  Retained uncounted as the source-local context the proof needs: lines 862--875.  Nothing was omitted from any retained span, so the package carries no omission record.  No retained line contains a citation, so the wrapper carries no bibliography and no bibliography entry.  No retained line contains a figure environment, an image-inclusion command or drawing code, so no image is delivered and the package declares no asset dependency.  Editorial formatting only, in addition: an AMS \texttt{amsart} preamble that restates the article's own declarations and macros used by the retained lines, namely the environments \texttt{lemma}, \texttt{proposition} on separate counters, together with the standard packages those lines need.  The retained environments print in this document as Lemma 1, Proposition 1 in order of appearance, whereas the article prints the same items with the numbering of its own full document.  The complete native source of the article as distributed in the arXiv e-print for this exact version is bundled unchanged as \texttt{sources/209118/source0.tex}.
\begin{lemma}
\label{lem:torsion-root-count}
Let $\phi$ and $\psi$ be Drinfeld modules of the same rank over an $A$-field,
and, for $i=1,\ldots,\ell$, let $\mathfrak q_i$ be distinct monic
irreducible polynomials of degree $m$, prime to the $A$-characteristic. For
every nonzero $u\in\Hom(\phi,\psi)$,
$$
\sum_{i=1}^{\ell}
\dim_{\vF_{q^m}}
\ker\left(u|_{\phi[\mathfrak q_i]}\right)
\leq
\left\lfloor\frac{\deg_\tau(u)}{m}\right\rfloor.
$$
\end{lemma}

\begin{proposition}
For the code $\mathcal C_{s,f}$, one has
$$
\dim_{\vF_q}\mathcal C_{s,f}=mrt-c
$$
and
$$
d_R(\mathcal C_{s,f})=r-t+1.
$$
In particular, $\mathcal C_{s,f}$ is exactly $c$ below the additive Singleton bound for its minimum rank distance.
\end{proposition}

\begin{proof}
For every nonzero $u\in M_s$, Lemma~\ref{lem:torsion-root-count}, applied
with $\ell=1$ and $\mathfrak q_1=f$, gives
$$
\dim_{\vF_{q^m}}\ker\left(u|_{\phi[f]}\right)
\leq
\left\lfloor\frac{tm-1}{m}\right\rfloor
=t-1.
$$
Therefore
$$
\rk_{\vF_{q^m}}\rho_f(u)
=
r-\dim_{\vF_{q^m}}\ker\left(u|_{\phi[f]}\right)
\geq r-t+1.
$$
Hence
$$
d_R(\mathcal C_{s,f})\geq r-t+1.
$$

We now compute the dimension. By the stabilization theorem,
$$
\dim_{\vF_q}M_s
=
r(s+1)-c
=
rtm-c.
$$
Moreover, the restriction map $\rho_f$ is injective on $M_s$. Indeed, if $0\neq u\in M_s$ were in the kernel of $\rho_f$, then $u|_{\phi[f]}=0$, so
$$
\dim_{\vF_{q^m}}\ker\left(u|_{\phi[f]}\right)=r.
$$
This contradicts the preceding bound, since $t<r$. Therefore $\rho_f$ is injective on $M_s$, and consequently
$$
\dim_{\vF_q}\mathcal C_{s,f}
=
\dim_{\vF_q}M_s
=
mrt-c.
$$

It remains to show that the lower bound on the minimum distance is sharp. Suppose, for a contradiction, that
$$
d_R(\mathcal C_{s,f})\geq r-t+2.
$$
Then the additive Singleton bound for rank-metric codes in $\Mat_{r\times r}(\vF_{q^m})$ gives
$$
\dim_{\vF_q}\mathcal C_{s,f}
\leq
mr\bigl(r-(r-t+2)+1\bigr)
=
mr(t-1).
$$
On the other hand, we have proved that
$$
\dim_{\vF_q}\mathcal C_{s,f}=mrt-c.
$$
The stabilization-range hypothesis gives
$$
tm-1\geq rc,
$$
and hence
$$
c\leq \frac{tm-1}{r}<\frac{tm}{r}.
$$
Since $t<r$, we have $c<m$, and in particular $c<mr$. Therefore
$$
mrt-c>mrt-mr=mr(t-1),
$$
contradicting the Singleton bound above. Hence
$$
d_R(\mathcal C_{s,f})=r-t+1.
$$

For this value of the minimum distance, the additive Singleton bound is
$$
\dim_{\vF_q}\mathcal C\leq mr\bigl(r-(r-t+1)+1\bigr)=mrt.
$$
Since
$$
\dim_{\vF_q}\mathcal C_{s,f}=mrt-c,
$$
the code $\mathcal C_{s,f}$ is exactly $c$ below the additive Singleton bound for its actual minimum rank distance.
\end{proof}

\end{document}
